\chapter{Статический анализ}
В процессе генерации условий корректности "в лоб" может порождаться большое число условий корректности соответствующих недостижимым путям в программе. Для того что бы этого избежать используется система статического анализа.

Сначала определим поля глобального и локального контекста.

Глобальный контекст в данном случае не нужен.

Поля локального контекста:

\begin{itemize}
    \item \textbf{currentProcess} -- имя текущего процесса.
    \item \textbf{processesToStart} -- список процессов для запуска.
    \item \textbf{collector} -- последовательность словарей с атрибутами
\end{itemize}

Словарь с атрибутами имеет следующие поля:
\begin{itemize}
    \item \textbf{processChange} -- словарь, отображающий имя процесса в c$\in$\{"start","stop","error"\} и обозначающий какое дествия было произведено с процессом после исполнения указанного оператора.
    \item \textbf{potProcessChange} -- список словарей с полями "process" $\rightarrow$ имя процесса и "change" $\rightarrow$ c$\in$\{"start","stop","error"\}, обозначающий какое действие мгло быть произведено с процессом после исполнения указанного оператора.
    \item \textbf{reset} -- булево значение, обозначающее был ли обновлен таймер процесса.
    \item \textbf{stateChanged} -- булево значение, обозначающее было ли изменено текущее состояние процесса.
    \item \textbf{changesTo} -- список имен процессных состояний, обозначающих в какое состояние может перейти теущий процесс, по результату исполнения оператора.
\end{itemize}

\section{Вспомогательные функции}
\begin{lstlisting}[style=PseudoStyle,mathescape=true]
fun addConsAttributes(oldAttr,newAttr){
    procChange = oldAttr.processChange
    potProcChange = oldAttr.potProcessChange
    newProcChange = newAttr.processChange
    newPotProcChange = newAttr.potProcessChange
    procs = ((definedProcChange(procChange) $\cap$ definedPotProcChange(newAttr.potProcessChange)) $\setminus$ definedProcChange(newAttr.potProcessChange))
    newChangesTo = newAttr.changesTo
    if(newAttr.reset){
        oldAttr.reset = true
    }
    if(newAttr.stateChanged){
        oldAttr.stateChanged = true
    }
    if(newChangesTo){
        if(member(null,newChangesTo)){
            newChangesTo = newChangesTo $\cup$ oldAttr.changesTo
        }
    }else{
        newChangesTo = oldAttr.changesTo
    }
    for proc in procs{
        potStart = memberIf(function(el){
            return el.process == proc && el.change == "start"
        }, potProcChange)
        potStop = memberIf(function(el){
            return el.process == proc && el.change == "stop"
        },potProcChange)
        potError = memberIf(function(el){
            return el.process == proc && el.change == "error"
        },potProcChange)
        if(procChange[proc] == "start" && (potStop || potError)){
            procChange[proc] = null
        }
        if(procChange[proc] == "stop" && (potStart || potError)){
            procChange[proc] = null
        }
        if(procChange[proc] == "error" && (potStart || potStop)){
            procChange[proc] = null
        }
    }
    for proc in newProcChange.keys {
        procChange[proc] = newProcChange[proc]
    }
    ppc1 = removeIf(function(change){
        return member(change.process,definedProcChange(newProcChange))
    }, potProcChange)
    ppc2 = union(ppc1,newPotProcChange)
    oldAttr.processChange = procChange
    oldAttr.potProcessChange = ppc2
    return oldAttr
}
fun consAttributes(attrList){
    if(attrList == null || attrList.length == 0){
        return createAttributes({})
    }
    result = clone(attrList[0])
    for i = 1; i < attrList.length; i++{
        result = addConsAttributes(result,attrList[i])
    }
    return result
}

fun addParAttributes(oldAttr,newAttr){
    procChange = oldAttr.processChange
    potProcChange = oldAttr.potProcessChange
    newProcChange = newAttr.processChange
    newPotProcChange = newAttr.potProcessChange
    newChangesTo = newAttr.changesTo
    if(newAttr.reset == null){
        newAttr.reset = null
    }
    if(newAttr.stateChanged == null){
        newAttr.stateChanged = null
    }
    if(newChangesTo){
        if (oldAttr.changesTo){
            newChangesTo = (newChangesTo $\cup$ oldAttr.changesTo)
        }else{
            newChangesTo = (newChangesTo $\cup$ [null])
        }
    }else{
        if (oldAttr.changesTo){
            newChangesTo = (oldAttr.changesTo $\cup$ [null])
        }else{
            null
        }
    }
    for proc in definedProcChange(procChange){
        if(!member(proc,definedProcChange(newProcChange)) ||
           newProcChange[proc] != procChange[proc]){
            procChange[proc] = null
        }
    }
    ppc1 = (potProcChange $\cup$ newPotProcChange)
    oldAttr.processChange = procChange
    oldAttr.potProcessChange = ppc1
    return oldAttr
}
fun addParAttributesConc(oldAttr){
    procChange = oldAttr.processChange
    potProcChange = oldAttr.potProcessChange
    definedProc = definedProcChange(procChange)
    ppc = []
    for elem in potProcChange{
        procName = elem.process
        if(member(procName,definedProc)){
            candidate = {process:procName,change:procChange[procName]}
            if(!exists(function(left){
                return left.process == candidate.process &&
                       left.change == candidate.change
            },ppc)){
                ppc.pushback(candidate)
            }
        }else{
            ppc.pushback(elem)
        }
    }
    oldAttr.potProcessChange = ppc
    return oldAttr
}
fun parAttributes(attrList){
    result = clone(attrList[0])
    for i = 1; i < attrList.length; i++{
        result = addParAttributes(result,attrList[i])
    }
    return addParAttributesConc(result)
}
\end{lstlisting}

\section{Функция расстановки базовых аттрибутов}
Определим функцию attributePrepare(st) расставляющую аттрибуты по операторам.

\begin{lstlisting}[style=PseudoStyle,mathescape=true]
fun attributePrepare(resetTimer) {
    attr = createAttributes({
        reset: true
    })
    resetTimer.analysisAttributes = attr
    return attr
}

fun attributePrepare(setState) {
    state = setState.state
    chanesTo = [state]
    attr = createAttributes({
        reset: true,
        chanesTo: chanesTo
    })
    setState.analysisAttributes = attr
    return attr 
}

fun attributePrepare(restartProcess) {
    curProc = ctx.currentProcess
    curState = ctx.currentState
    fstate = firstState(env, curProc)
    if (fstate != curState) {
        chanesTo = [fstate]
    } else {
        chanesTo = null
    }
    attr = createAttributes({
        processChange: {curProc:"start"},
        potProcessChange: [(curProc,"start")],
        reset: true,
        chanesTo: chanesTo,
        stateChanged: fstate != curState
    })
    restartProcess.analysisAttributes = attr
    return attr
}

fun attributePrepare(startProcess) {
    process = startProcess.process
    curProc = ctx.currentProcess
    if (process == curProc) {
        return attributePrepare(createRestartProcess())
    }
    attr = createAttributes({
        processChange: {process:"start"},
        potProcessChange: [(process,"start")]
    })
    startProcess.analysisAttributes = attr
    return attr
}

fun attributePrepare(stopCurrentProcess) {
    curProc = ctxcurrentProcess
    attr = createAttributes({
        processChange: {curProc:"stop"},
        potProcessChange: [{process:CurProc,change:"stop"}],
        reset: true,
        chanesTo: [ "stop" ],
        stateChanged: true
    })
    stopCurrentProcess.analysisAttributes = attr
    return attr
}

fun attributePrepare(stopProcess){
    process = stopProcess.process
    curProc = ctx.currentProcess
    if(process == curProc){
        attr = attributePrepare(createStopCurrentProcess())
        stopProcess.analysisAttributes = attr
        return attr
    }
    attr =createAttributes({
        processChange:{process:"stop"},
        potProcessChange:[{process:process,change:"stop"}]
    })
    stopProcess.analysisAttributes = attr
    return attr
}

fun attributePrepare(errorCurrentProcess){
    curProc = ctx.currentProcess
    attr = createAttributes({
        processChange:{curProc:"error"},
        potProcessChange:[{process:curProc,change:"error"}],
        reset:true,
        chanesTo:["error"],
        stateChanged:true
    })
    errorCurrentProcess.analysisAttributes = attr
    return attr
}

fun attributePrepare(errorProcess){
    process = errorProcess.process
    curProc = ctx.currentProcess
    if(process == curProc){
        attr = attributePrepare(createErrorCurrentProcess())
        errorProcess.analysisAttributes = attr
        return attr
    }
    attr = createAttributes({
        processChange:{process:"error"},
        potProcessChange:[{process:process,change:"error"}]
    })
    errorProcess.analysisAttributes = attr
    return attr
}

fun attributePrepare(statementList){
    collectedAttr = createAttributes({})
    for st in statementList {
        val = attributePrepare(st)
        addConsAttributes(collectedAttr,val)
    }
    return collectedAttr
}
fun attributePrepare(timeoutStatement){
    val = attributePrepare(timeoutStatement.statements)
    res = addParAttributesConc(
        addParAttributes(createAttributes({}),val)
    )
    timeoutStatement.analysisAttributes = res
    return res
}
fun attributePrepare(ifThen){
    val = attributePrepare(ifThen.then)
    res = addParAttributesConc(
        addParAttributes(createAttributes({}),val)
    )
    ifThen.analysisAttributes = res
    return res
}
fun attributePrepare(ifThenElse){
    thenVal = attributePrepare(ifThenElse.then)
    elseVal = attributePrepare(ifThenElse.else)
    res = addParAttributesConc(
        addParAttributes(thenVal,elseVal)
    )
    ifThenElse.analysisAttributes = res
    return res
}

fun attributePrepare(switchS){
    css = switchStatement.cases
    collected = []
    for cs in css {
        val = attributePrepare(cs)
        cs.analysisAttributes = val
        collected.pushback(val)
    }
    def = switchStatement.default
    if(def){
        defVal = attributePrepare(def)
    }else{
        defVal = createAttributes({})
    }
    collected.pushback(defVal)
    parAttr = parAttributes(collected)
    switchStatement.analysisAttributes = parAttr
    return parAttr
}

fun attributePrepare(defaultStatement){
    attr = attributePrepare(defaultStatement.statements)
    defaultStatement.analysisAttributes = attr
    return attr
}
fun attributePrepare(caseStatement){
    attr = attributePrepare(caseStatement.statements)
    caseStatement.analysisAttributes = attr
    return attr
}
fun attributePrepare(statementBlock){
    attr = attributePrepare(statementBlock.statements)
    statementBlock.analysisAttributes = attr
    return attr
}
fun attributePrepare(forStatement){
    val = attributePrepare(forStatement.statement)
    res =addParAttributesConc(   
        addParAttributes(createAttributes({}),val)
    )
    forStatement.analysisAttributes = res
    return res
}

fun attributePrepare(stateDeclaration){
    val = attributePrepare(stateDeclaration.statements)
    stateDeclaration.analysisAttributes = val
    return val
}
fun attributePrepare(processDeclaration){
    states = processDeclaration.states
    collected = []
    for st in states{
        val = attributePrepare(st)
        collected.pushback(val)
    }
    procAttr = parAttributes(collected)
    processDeclaration.analysisAttributes = procAttr
    return procAttr
}
fun attributePrepare(programDeclaration){
    procs = programDeclaration.processes
    for p in procs{
        attributePrepare(p)
    }
}

\end{lstlisting}

\section{Расстановка дополнительных атрибутов}
Для дальнейшего анализа используются несколько дополнительных атрибутов:
\begin{itemize}
    \item \textbf{reachS} -- обозначет, может ли текущий процесс быть переведен в состояние "stop".
    \item \textbf{startS} -- обозначает, начинается ли текщий процесс в состоянии "stop".
    \item \textbf{reachE}  -- обозначет, может ли текущий процесс быть переведен в состояние "error".
    \item \textbf{group} -- целочисленное значение, определяющее группу процесса.
\end{itemize}

Если несколько процессов обьеденины в одну группу, то они запускаются, останавливаются и останваливаются с ошибкой одновременно.

Означивание атрибутов \textbf{reachS}, \textbf{startS} и \textbf{reachE} определяется следующим алгоритмом:

\begin{lstlisting}[style=PseudoStyle,mathescape=true]
fun setReachS(procs){
    procsSetStop = []
    for p in procs {
        for change in p.potProcChange {
            if change.change == "stop" {
                procsSetStop = procsSetStop $\cup$ change.process}}}
    for p in procs {
        if "stop" $\in$ p.changesTo  || p.name $\in$ procsSetStop{
            p.analysisAttributes.reachS = true}}
}

fun setReachE(procs){
    procsSetStop = []
    for p in procs {
        for change in p.potProcChange {
            if change.change == "error" {
                procsSetStop = procsSetStop $\cup$ change.process}}}
    for p in procs {
        if "error" $\in$ p.changesTo  || p.name $\in$ procsSetStop{
            p.analysisAttributes.reachS = true}}
}

// В данном определении исполтзуется функция tailAfter (list,el) -- она возвращает список элементов list идущих после el.
fun setStartS(procs){
    for proc in procs {
        if proc.active {
            proc.analysisAttributes.startS = false
        } else {
            proc.analysisAttributes.startS = true
        }
    }
    revProcs = reverse(procs)
    for curProc in procs{
        curName = curProc.name
        tailRev = tailAfter(revProcs,curProc)
        headList = reverse(tailRev)
        found = findIf(function(prevProc){
            prevStartS = prevProc.analysisAttributes.startS
            firstState = prevProc.states[0]
            procChange = firstState.analysisAttributes.processChange[curName]
            if(prevStartS && procChange == "start"){
                midList = tailAfter(prevProc,headList)
                blocked = findIf(function(midProc){
                    return pairInPotProcChange(
                        curName,
                        "stop",
                        midProc.analysisAttributes.potProcessChange
                    )
                }, midList)
                return !blocked
            }
            return false
        }, headList)
        if(found){
            curProc.analysisAttributes.startS = false
        }
    }
} 
\end{lstlisting}

Определение атрибута \textbf{group} определяется следующими функциями:

\begin{lstlisting}[style=PseudoStyle,mathescape=true]
fun setsInter(setsSet,procSet){
    coll = []
    for el in setsSet {
        inter = el $\cap$ procSet 
        diff = el $\setminus$ procSet 
        tmp = inter $\cup$ coll
        coll = diff $\cup$ tmp
    }
    return coll
}

fun procId(procName){
    return env.processId[procName]
}

fun setsDivCore(setsSet,hpc,procChange){
    currentProcess = ctx.currentProcess
    currentState = ctx.currentState
    nhpc = clone(procChange) $\cup$ hpc
    startPpred = []
    startPsucc = []
    stopPpred = []
    stopPsucc = []
    errorPpred = []
    errorPsucc = []
    curPid = procId(env,currentProcess)
    for proc in nhpc.keys {
        pid = procId(env,proc)
        change = nhpc[proc]
        if(change == 'start'){
            if(pid < curPid){
                startPpred.pushback(proc)
            }
            else if(
                pid > curPid ||
                (pid == curPid && currentState == firstState(env,currentProcess))
            ){
                startPsucc.pushback(proc)
            }
        }
        if(change == 'stop'){
            if(pid < curPid){
                stopPpred.pushback(proc)
            }
            else if(pid > curPid){
                stopPsucc.pushback(proc)
            }
        }
        if(change == 'error'){
            if(pid < curPid){
                errorPpred.pushback(proc)
            }
            else if(pid > curPid){
                errorPsucc.pushback(proc)
            }
        }
    }
    newSets = setsSet
    groups = [startPpred,
        startPsucc,
        stopPpred,
        stopPsucc,
        errorPpred,
        errorPsucc]

    for grp in groups {
        newSets = setsInter(newSets,grp)
    }

    return [nhpc,newSets]
}

fun groupStaticAnalysisPrepare(program){
    program = closure.instance
    procs = program.processes
    nonStart = []
    start = []
    for proc in procs {
        proc.analysisAttributes.group = uid(proc)
        if(!proc.analysisAttributes.startS){
            nonStart.pushback(proc.name)
        } else {
            start.pushback(proc.name)
        }
    }
    setsDiv(program, [nonStart,start],{})
    setsSet = setsDiv(program, [nonStart,start],{})
    index = 0
    for group in setsSet {
        for gel in group {
            for proc in procs {
                if(gel == proc.name){
                    proc.analysisAttributes.group = index
                }
            }
        }
        index = index + 1
    }
}

fun setsDiv(resetTimer,setsSet,hpc){
    procChange = resetTimer.analysisAttributes.procChange
    return setsDivCore(setsSet,hpc,procChange)
}
fun setsDiv(setState,setsSet,hpc){
    procChange = setState.analysisAttributes.procChange
    return setsDivCore(setsSet,hpc,procChange)
}
fun setsDiv(restartProcess,setsSet,hpc){
    procChange = restartProcess.analysisAttributes.procChange
    return setsDivCore(setsSet,hpc,procChange)
}
fun setsDiv(startProcess,setsSet,hpc){
    procChange = startProcess.analysisAttributes.procChange
    return setsDivCore(setsSet,hpc,procChange)
}
fun setsDiv(stopCurrentProcess,setsSet,hpc){
    procChange = stopCurrentProcess.analysisAttributes.procChange
    return setsDivCore(setsSet,hpc,procChange)
}
fun setsDiv(stopProcess,setsSet,hpc){
    procChange = stopProcess.analysisAttributes.procChange
    return setsDivCore(setsSet,hpc,procChange)
}
fun setsDiv(errorCurrentProcess,setsSet,hpc){
    procChange = errorCurrentProcess.analysisAttributes.procChange
    return setsDivCore(setsSet,hpc,procChange)
}
fun setsDiv(errorProcess,setsSet,hpc){
    procChange = errorProcess.analysisAttributes.procChange
    return setsDivCore(setsSet,hpc,procChange)
}
fun setsDiv(statementList,setsSet,hpc){
    rst = statementList.slice(1)
    val = setsDiv(statementList[0],setsSet,hpc)
    for st in rst{
        val = setsDiv(st,val[1],val[0])
    }
    return val
}
fun setsDiv(timeoutStatement,setsSet,hpc){
    procChange = timeoutStatement.analysisAttributes.procChange
    newHpcSets = setsDivCore(setsSet,hpc,procChange)
    return setsDiv(
        timeoutStatement.statements,
        newHpcSets[1],
        newHpcSets[0]
    )
}

fun setsDiv(ifThenStatement,setsSet,hpc){
    procChange = ifThenStatement.analysisAttributes.procChange
    newHpcSets = setsDivCore(setsSet,hpc,procChange)
    return setsDiv(
        ifThenStatement.then,
        newHpcSets[1],
        newHpcSets[0]
    )
}
fun setsDiv(ifThenElseStatement,setsSet,hpc){
    procChange = ifThenElseStatement.analysisAttributes.procChange
    newHpcSets = setsDivCore(setsSet,hpc,procChange)
    thenVal = setsDiv(
        ifThenElseStatement.then,
        newHpcSets[1],
        newHpcSets[0]
    )
    return setsDiv(
        ifThenElseStatement.else,
        thenVal[1],
        newHpcSets[0]
    )
}
fun setsDiv(switchStatement,setsSet,hpc){
    procChange = switchStatement.analysisAttributes.procChange
    newHpcSets = setsDivCore(setsSet,hpc,procChange)
    val = setsDiv(
        switchStatement.cases[0],
        newHpcSets[1],
        newHpcSets[0]
    )
    for cs in switchStatement.cases.slice(1){
        val = setsDiv(cs,val[1],newHpcSets[0])
    }
    if(switchStatement.default){
        return setsDiv(
            switchStatement.default,
            val[1],
            newHpcSets[0]
        )
    }
    return val
}
fun setsDiv(caseStatement,setsSet,hpc){
    procChange = caseStatement.analysisAttributes.procChange
    newHpcSets = setsDivCore(setsSet,hpc,procChange)
    return setsDiv(
        caseStatement.statements,
        newHpcSets[1],
        newHpcSets[0]
    )
}
fun setsDiv(defaultStatement,setsSet,hpc){
    procChange = defaultStatement.analysisAttributes.procChange
    newHpcSets = setsDivCore(setsSet,hpc,procChange)
    return setsDiv(
        defaultStatement.statements,
        newHpcSets[1],
        newHpcSets[0]
    )
}
fun setsDiv(statementBlock,setsSet,hpc){
    procChange = statementBlock.analysisAttributes.procChange
    newHpcSets = setsDivCore(setsSet,hpc,procChange)
    return setsDiv(
        statementBlock.statements,
        newHpcSets[1],
        newHpcSets[0]
    )
}
fun setsDiv(expressionStatement,setsSet,hpc){
    return setsSet
}
fun setsDiv(stateDeclaration,setsSet,hpc){
    ctx.currentState = stateDeclaration.name
    procChange = stateDeclaration.analysisAttributes.procChange
    newHpcSets = setsDivCore(setsSet,hpc,procChange)
    return setsDiv(
        stateDeclaration.statements,
        newHpcSets[1],
        newHpcSets[0]
    )
}
fun setsDiv(processDeclaration,setsSet,hpc){
    ctx.currentProcess = processDeclaration.name
    procChange = processDeclaration.analysisAttributes.procChange
    newHpcSets = setsDivCore(setsSet,hpc,procChange)
    val = setsDiv(
        processDeclaration.states[0],
        newHpcSets[1],
        newHpcSets[0]
    )
    for st in processDeclaration.states.slice(1){
        val = setsDiv(st,val[1],newHpcSets[0])
    }
    return val
}
fun setsDiv(programDeclaration,setsSet,hpc){
    procChange = programDeclaration.analysisAttributes.procChange
    newHpcSets = setsDivCore(setsSet,hpc,procChange)
    val = setsDiv(
        programDeclaration.processes[0],
        newHpcSets[1],
        newHpcSets[0]
    )
    for p in programDeclaration.processes.slice(1){
        val = setsDiv(p,val[1],newHpcSets[0])
    }
    return val
}

\end{lstlisting}

\section{Применение атрибутов}

Для использования атрибутов применяются наборы правил. Некоторые из них опираются на поле локального контекста curAttr. Данное поле хранит атрибуты статического анализа полученные на пройденом пути. 
Для применения общих правил:

\begin{lstlisting}[style=PseudoStyle,mathescape=true]
fun isProcessActivity(term,process,activity){
    if(isInstance(term,"processActivity")){
        return term.process == process &&
               term.activity == activity
    }
    if(isInstance(term,"processActivityBlock")){
        for act in term {
            if(isProcessActivity(act,process,activity)){ return true }
        }
    }
    return false
}

fun simpleStaticAnalysisWork(term){
    currentProcess = ctx.currentProcess
    curFirstState = firstState(env,currentProcess)
    curCond = reverse(ctx.curCond)
    if(isInstance(term,"pstateCompare")){
        if(checkRule1(term)){ return false }
        if(checkRule2(term)){ return false }
        if(checkRule3(term,curCond,curFirstState,currentProcess)){ return false }
        if(checkRule4(term,curCond,currentProcess)){ return false }
        if(checkRule5(term,curCond,currentProcess)){ return false }
        if(checkRule6(term,curCond,currentProcess)){ return false }
        return true
    }
    if(isInstance(term,"ltimeCheck")){
        if(checkRule7(term,curCond,currentProcess)){ return false }
        return true
    }
    if(isInstance(term,"processActivity")){
        if(checkRule8(term,curCond)){ return false }
        if(checkRule9(term,curCond)){ return false }
        return true
    }
    if(isInstance(term,"processActivityBlock")){
        for formula in term {
            if(checkRule8(term,curCond)){ return false }
            if(checkRule9(term,curCond)){ return false }
        }
        return true
    }
    return true
}

fun checkRule1(term){
    return term.pstate == "error" &&
           !ctx.curAttr.reachE
}

fun checkRule2(term){
    return term.pstate == "stop" &&
           !ctx.curAttr.reachS &&
           !ctx.curAttr.startS
}

fun checkRule3(term,curCond,curFirstState,currentProcess){
    return term.pstate != curFirstState &&
           ctx.curAttr.processChange[currentProcess] == 'start'
}

fun checkRule4(term,curCond,currentProcess){
    return term.pstate != "stop" &&
           ctx.curAttr.processChange[currentProcess] == 'stop'
}

fun checkRule5(term,curCond,currentProcess){
    return term.pstate != "error" &&
           ctx.curAttr.processChange[currentProcess] == 'error'
}

fun checkRule6(term,curCond,currentProcess){
    for formula in curCond {
        p = term.pstate
        if((p == "stop" || p == "error") &&
           isProcessActivity(formula,currentProcess,'active')){ return true }
        if(p != "stop" && p != "error" &&
           isProcessActivity(formula,currentProcess,'inactive')){ return true }
        if(p != "stop" &&
           isProcessActivity(formula,currentProcess,'stop')){ return true }
        if(p != "error" &&
           isProcessActivity(formula,currentProcess,'error')){ return true }
        if(p == "stop" &&
           isProcessActivity(formula,currentProcess,'nonstop')){ return true }
        if(p != "error" &&
           isProcessActivity(formula,currentProcess,'nonerror')){ return true }
    }
    return false
}

fun checkRule7(term,curCond,currentProcess){
    return !ctx.curAttr.reset
}

fun checkRule8(term,curCond){
    proc = term.process
    act = term.activity
    curChange = ctx.curAttr.processChange[proc]
    // sub1
    if(act == 'active'){
        found = false
        for formula in curCond {
            if(isProcessActivity(formula,proc,'inactive') ||
               isProcessActivity(formula,proc,'stop') ||
               isProcessActivity(formula,proc,'error')){
                found = true
            }
        }
        if(found && curChange != 'start'){ return true }
    }
    // sub2
    if(act == 'inactive'){
        found = false
        for formula in curCond {
            if(isProcessActivity(formula,proc,'active')){ found = true }
        }

        if(found && curChange != 'stop' && curChange != 'error'){ return true }
    }
    // sub3
    if(act == 'stop'){
        found = false
        for formula in curCond {
            if(isProcessActivity(formula,proc,'active') ||
               isProcessActivity(formula,proc,'nonstop') ||
               isProcessActivity(formula,proc,'error')){
                found = true
            }
        }
        if(found && curChange != 'stop'){ return true }
    }
    // sub4
    if(act == 'error'){

        found = false
        for formula in curCond {
            if(isProcessActivity(formula,proc,'active') ||
               isProcessActivity(formula,proc,'nonerror') ||
               isProcessActivity(formula,proc,'stop')){
                found = true
            }
        }
        if(found && curChange != 'error'){ return true }
    }
    // sub5
    if(act == 'nonstop'){
        found = false
        for formula in curCond {
            if(isProcessActivity(formula,proc,'stop')){ found = true }
        }
        if(found && curChange != 'start' && curChange != 'error'){ return true }
    }
    // sub6
    if(act == 'nonerror'){
        found = false
        for formula in curCond {
            if(isProcessActivity(formula,proc,'error')){ found = true }
        }
        if(found && curChange != 'start' && curChange != 'stop'){ return true }
    }
    return false
}

fun checkRule9(term,curCond){
    proc = term.process
    act = term.activity
    curChange = ctx.curAttr.processChange[proc]
    if(act == 'active'){
        if(curChange != 'stop' && curChange != 'error'){ return true }
    }
    if(act == 'inactive'){
        found = false
        for formula in curCond {
            if(isProcessActivity(formula,proc,'active')){ found = true }
        }
        if(found && curChange != 'start'){ return true }
    }
    if(act == 'stop'){
        if(curChange != 'start' && curChange != 'error'){ return true }
    }
    if(act == 'error'){
        if(curChange != 'stop' && curChange != 'start'){ return true }
    }
    if(act == 'nonstop'){
        found = false
        for formula in curCond {
            if(isProcessActivity(formula,proc,'stop')){ found = true }
        }
        if(found && curChange != 'stop'){ return true }
    }
    if(act == 'nonerror'){
        found = false
        for formula in curCond {
            if(isProcessActivity(formula,proc,'error')){ found = true }
        }
        if(found && curChange != 'error'){ return true }
    }
    return false
}
\end{lstlisting}

Для применения правил группировки:

\begin{lstlisting}[style=PseudoStyle,mathescape=true]
fun groupStaticAnalysisWork(term){
    currentProcess = ctx.currentProcess
    revCurCond = ctx.curCond
    curCond = reverse(revCurCond)
    if(isInstance(term,"pstateCompare")){
        r1 = checkRule1(term,curCond,currentProcess)
        if(r1){ return false }
        r2 = checkRule2(term,curCond,currentProcess)
        if(r2){ return false }
        r3 = checkRule3(term,curCond,currentProcess)
        if(r3){ return false }
        r4 = checkRule4(term,curCond,currentProcess)
        if(r4){ return false }
        return true
    }
    return true
}
fun getProcGroup(procName){
    for proc in env.program.processes {
        if(proc.name == procName){
            return proc.analysisAttributes.group
        }
    }
    return null
}
fun getProc(procName){
    for proc in env.program.processes {
        if(proc.name == procName){
            return proc
        }
    }
    return null
}
fun checkRule1(term,curCond,currentProcess){
    for formula in curCond {
        if(isInstance(formula,"pstateCompare")){
            otherProcName = formula.process
            if(otherProcName != currentProcess &&
               getProcGroup(otherProcName) == getProcGroup(currentProcess)){
                termStop = (term.change == 'stop')
                formulaStop = (formula.change == 'stop')
                if(termStop != formulaStop){
                    return true
                }}}
    }
    return false
}
fun checkRule2(term,curCond,currentProcess){
    for formula in curCond {
        if(isInstance(formula,"pstateCompare")){
            otherProcName = formula.process
            if(otherProcName != currentProcess &&
               getProcGroup(otherProcName) == getProcGroup(currentProcess)){
                termErr = (term.change == 'error')
                formulaErr = (formula.change == 'error')
                if(termErr != formulaErr){
                    return true
                }}}
    }
    return false
}
fun checkRule3(term,curCond,currentProcess){
    curStateName = ctx.currentState
    curProc = getProc(currentProcess)
    curFirstState = curProc.states[0]
    for formula in curCond {
        if(isInstance(formula,"pstateCompare")){
            otherProcName = formula.process
            otherProc = getProc(otherProcName)
            otherFirstState = otherProc.states[0]
            if(otherProcName != currentProcess &&
               getProcGroup(otherProcName) == getProcGroup(currentProcess) &&
               curFirstState.analysisAttributes.reachFrom == null &&
               curFirstState.name == curStateName &&
               curFirstState.analysisAttributes.stateChanged &&
               otherFirstState.name == otherProcName){
                return true
            }}
    }
    return false
}
fun checkRule4(term,curCond,currentProcess){
    curStateName = ctx.currentState
    curProc = getProc(currentProcess)
    curFirstState = curProc.states[0]
    for formula in curCond {
        if(isInstance(formula,"pstateCompare")){
            otherProcName = formula.process
            otherProc = getProc(otherProcName)
            otherFirstState = otherProc.states[0]
            if(otherProcName != currentProcess &&
               getProcGroup(otherProcName) == getProcGroup(currentProcess) &&
               otherFirstState.analysisAttributes.reachFrom == null &&
               otherFirstState.name == formula.pstate &&
               otherFirstState.analysisAttributes.stateChanged &&
               curFirstState.name == currentProcess){

                return true
            }}
    }
    return false
}
\end{lstlisting}

Вызов проверки соответствия терма правилам производится функцией checkCompatability(term)

\begin{lstlisting}[style=PseudoStyle,mathescape=true]
fun checkCompatability(term){
    return simpleStaticAnalysisWork(term) || groupStaticAnalysisWork(term)
}
\end{lstlisting}
